% META: {"id":"1SPE-DERLOCAL-EX-041","chapitre":"1SPE-DERIVATION-LOCAL","type_objet":"exercice","capacites":["1SPE-DERIVATION-LOCAL-C2"],"capacites_codes":["C2"],"methodes":["M1"],"parcours":2,"competences":["calculer","modeliser"],"duree_min":15,"mode_creation":"ex_nihilo","sources_inspiration":[],"fichier_tex":"chapitres/1SPE-DERIVATION-LOCAL/exercices/1SPE-DERLOCAL-EX-041.tex","corrige_tex":"chapitres/1SPE-DERIVATION-LOCAL/corriges/1SPE-DERLOCAL-CO-041.tex","status":"approved","origin":"generated","status_history":["generated","audited","accepted","approved"],"release_acceptance":"RELEASE_OWNER_BATCH_ACCEPTANCE","acceptance_closure_digest":"sha256:9b3ccf9a81c5520fbb7e03b7b2d3e7bf2057a02834b6d908bf3be5f49f3b553f","acceptance_receipt_digest":"sha256:4fad86f0282e0fa4e0419cca1ad6379ae7af5a280c04a4a90880f90a1c044242"}
% BEGIN-VERIFY
% from sympy import symbols, limit, expand, simplify
% t, h = symbols('t h')
% P = 500 + 100*t + 5*t**2
% assert P.subs(t, 4) == 980
% taux = (P.subs(t, 4+h) - P.subs(t, 4)) / h
% taux_expanded = expand(taux)
% assert limit(taux, h, 0) == 140
% END-VERIFY
\begin{exercice}{1SPE-DERLOCAL-EX-041}{2}{15}
La population d'une ville est modélisée par la fonction :
\[
P(t) = 500 + 100t + 5t^2
\]
où $t$ est le nombre d'années écoulées depuis 2010 et $P(t)$ est exprimé en milliers d'habitants.
\begin{enumerate}
\item Calculer $P(4)$ et exprimer $P(4+h)$ en fonction de $h$.
\item En déduire le taux de variation $\dfrac{P(4+h) - P(4)}{h}$ et le simplifier.
\item Calculer la limite de ce taux lorsque $h$ tend vers $0$. Que représente ce nombre ?
\item Interpréter le résultat dans le contexte démographique.
\end{enumerate}
\end{exercice}
